\chapter{Giving the model tools}
\label{ch:tools}

The fastest way to make \nodwyn{} better at \emph{your} project is to give
the model tools that fit it: run the tests the way this project runs them,
read a design record by number, summarise the build output. A tool is a
Scheme procedure with a description, registered from your folder's
\file{.nodwyn/init.scm}.

\section{\texttt{init.scm}}

When a session starts in a folder, \nodwyn{} loads its prelude, then the
default loop, then \file{<folder>/.nodwyn/init.scm} if it exists.
Everything that file defines lives in the session's global environment, and
every tool it registers is offered to the model.

A good \code{init.scm} registers and defines; it does not \emph{do}. The
\nodwyn{} repository's own is five lines:

\begin{scm}
;; Loaded after the prelude when nodwyn starts in this repository (ADR 14).
;; Everything here registers tools; nothing here performs an action.
(load "dev/go.scm")
(load "dev/git.scm")
(load "dev/repo.scm")
\end{scm}

\code{load} evaluates a file's forms. It is a host primitive, so it goes
through the guard like a read, and a relative path starts at the folder.

\begin{careful}
\file{.nodwyn/} is a protected path: every write to it is a security
confirm, so no standing approval, turn approval or \code{-{}-yolo} can let
the model change your \code{init.scm} behind your back. Files you
\code{load} from elsewhere (\file{dev/} above) are ordinary files; if you
keep tools there, review changes to them as you would review code, because
they run with your authority in the next session.
\end{careful}

\section{\texttt{register-tool}}

\begin{scm}
(register-tool "plant-list"
  (lambda () (lines (fs-read "plants.txt")))
  '((doc . "List the plants growing in the garden.")
    (params . ())))
\end{scm}

Three arguments: a name, a procedure, and metadata. The rules:

\begin{itemize}
  \item \code{doc} and \code{params} are required. The doc is the model's
    only instruction about the tool, so write it for the model: what it
    does, what it returns, when to prefer it.
  \item Each parameter is \code{(name type doc)}. The model sends JSON; the
    arguments are parsed as Chapter~\ref{ch:strings} described and passed
    in the order you declared.
  \item The name must be new. A builtin cannot be shadowed, and a tool
    cannot be registered twice: the tool list is fixed for the whole
    session, so the model's view of its tools never shifts under it (and
    the model server can keep the start of the prompt cached).
  \item \code{irreversible}, \code{egress} and \code{variadic} default to
    \code{\#f}.
\end{itemize}

\begin{txt}
you> :eval (register-tool "plant-list" (lambda () 1) '((doc . "x") (params . ())))
error: register-tool: scheme: primitive "plant-list" already registered
you> :eval (register-tool "fs-read" (lambda () 1) '((doc . "x") (params . ())))
error: register-tool: scheme: primitive "fs-read" already registered
you> :eval (register-tool "nodoc" (lambda () 1) '((params . ())))
error: register-tool: doc and params are required
\end{txt}

\section{A garden toolkit}

Put this in \file{.nodwyn/garden.scm} and load it from \code{init.scm}:

\begin{scm}
;; Garden tools: named, documented wrappers the model can call.

(define (lines text)
  (filter (lambda (l) (not (string-empty? (string-trim l))))
          (string-split text "\n")))

(register-tool "plant-list"
  (lambda () (lines (fs-read "plants.txt")))
  '((doc . "List the plants growing in the garden.")
    (params . ())))

(register-tool "plant-table"
  (lambda (plants)
    (md-table '("plant" "letters")
              (map (lambda (p) (list p (string-length p))) plants)))
  '((doc . "Show plant names as a Markdown table with their lengths.")
    (params . (("plants" "list" "plant names, e.g. [\"basil\", \"mint\"]")))))

(register-tool "count-lines"
  (lambda (path)
    (string-trim (assoc-ref (shell-exec "wc" (list "-l" path)) "stdout")))
  '((doc . "Count the lines of a file in the garden.")
    (params . (("path" "string" "a file relative to the garden")))))
\end{scm}

(\code{md-table} is the one from Chapter~\ref{ch:playground}.) The session
banner now lists them, and \code{:tools} shows what the model reads:

\begin{txt}
tools count-lines fs-edit fs-find fs-list fs-read fs-write index-grep
      plant-list plant-table read-definition second-look shell-exec
      shell-read shell-send shell-start shell-stop

you> :eval (display (plant-table (plant-list)))
| plant | letters |
| --- | --- |
| tomatoes | 8 |
| basil | 5 |
you> :eval (count-lines "plants.txt")
⚑ guard confirm shell-exec: shell-exec is irreversible → yes
"2 plants.txt"
\end{txt}

Notice the last one. \code{count-lines} was registered as reversible, and
it still asked, because the guard sees \emph{both} the call to your tool
and the \code{shell-exec} inside it. A tool is never a way around the
guard; it is a way to give an action a good name and a good description.
Pair a tool that runs a command with a standing approval for exactly that
command (\code{commands = ["wc -l"]} in \file{.nodwyn/allow.toml}), and it
will run without asking.

\section{The pattern that wins: \texttt{run}}

\nodwyn's own \file{dev/go.scm} shows the pattern for command tools. One
helper renders a command's result the way a person would want to read it:

\begin{scm}
;; run: execute and render for the model: exit status, then stdout and
;; stderr as they were produced.
(define (run cmd args)
  (let ((r (shell-exec cmd args)))
    (format "exit ~a~%~a~a"
            (assoc-ref r "status")
            (assoc-ref r "stdout")
            (assoc-ref r "stderr"))))

(register-tool "go-test"
  (lambda (pkg) (run "go" (list "test" "-race" pkg)))
  '((doc . "Run the Go tests with the race detector for a package pattern such as ./internal/log or ./... . Returns the exit status and output.")
    (params . (("pkg" "string" "package pattern, e.g. ./internal/log or ./...")))))
\end{scm}

And one tool does something no shell command can: compose file primitives
in Scheme, with no process at all:

\begin{scm}
;; read-adr: find docs/adr/NNNN-*.md by number and read it.
(define (pad4 n)
  (let ((s (number->string n)))
    (string-append (substring "0000" 0 (- 4 (string-length s))) s)))

(define (adr-file n)
  (let ((prefix (format "~a-" (pad4 n))))
    (let loop ((names (fs-list "docs/adr")))
      (cond ((null? names) #f)
            ((string-prefix? prefix (car names)) (string-append "docs/adr/" (car names)))
            (else (loop (cdr names)))))))

(register-tool "read-adr"
  (lambda (n)
    (let ((f (adr-file n)))
      (if f (fs-read f) (error "no such ADR" n))))
  '((doc . "Read architecture decision record number n from docs/adr.")
    (params . (("n" "integer" "the ADR number, e.g. 10")))))
\end{scm}

When the tool raises an error, the model reads the error text as the tool
result and can try again: \code{(error "no such ADR" n)} is a message to
the model, so make it useful.

\section{Good tools, bad tools}

\begin{tip}
\begin{itemize}
  \item \textbf{Name the project's way of doing things.} A model will guess
    \code{pytest}; your tool knows it is \code{tox -e py312 -{}- -x}.
  \item \textbf{Return text a model can use}: an exit status first, then the
    output; a table rather than raw data; the end of a long log, where the
    summary is.
  \item \textbf{Prefer list parameters to variadic ones.} The model's tool
    schema cannot express ``any number of arguments''; a \code{list}
    parameter it understands.
  \item \textbf{Keep tools read-only when you can}, and let the model use
    \code{fs-edit} for changes, so every change still shows you a diff.
  \item \textbf{Don't mirror every command.} Each tool adds to the prompt
    of every request. A handful of sharp tools beats thirty blunt ones.
\end{itemize}
\end{tip}

\section{Trying things without a model}

You do not need the model to test your tools. \code{:eval} calls them
directly, with your origin, through the same guard:

\begin{txt}
you> :eval (plant-list)
("tomatoes" "basil")
you> :eval (primitive-info "plant-table")
((doc . "Show plant names as a Markdown table with their lengths.")
 (params ("plants" "list" "plant names, e.g. [\"basil\", \"mint\"]"))
 (variadic . #f) (egress . #f) (irreversible . #f) (host . #f)
 (origin . "scheme"))
\end{txt}

\begin{puzzle}[ 13: a harvest tool]
Register a tool \code{harvest-report} with no parameters that reads
\file{harvest.csv} (lines like \code{tomato,12}) and returns a Markdown
table sorted by count, largest first, with a total row at the end.\ans{13}
\end{puzzle}
